\section{Conclusion}\label{sec:conclusion}

This paper introduced Green MGGP, a bi-objective multigene GP method for polynomial NARX models that minimises the
multiple-shooting simulation error and the arithmetic cost of the executed, duplicate-free polynomial with
Pareto-based selection, where the cost weights multiplications relative to additions by a platform parameter $\rho$,
and evaluated it on the regression examples of the \texttt{mggp} package over \nSeeds{} independent runs. Timing the
free-run simulation of the selected models showed that the execution time follows the number of operations, with a
multiplication costing about as much as an addition, so the unit weight was adopted. With it, the trade-off sets of
Green MGGP reached a higher hypervolume than those of the single-objective package and of its lexicographic
bi-objective mode in every run, both in identification and in free-run validation; its selected models needed
$17$--$32\%$ fewer operations than the single-objective ones, in terms of median cost, and $13$--$35\%$ less energy when
executed as canonical polynomials; and on the noisy bilinear system its trade-off set always contained the true
regressors. The same hypervolume advantage was obtained with the bit-level Karatsuba weight of the Green-Box framework,
although the two weights selected different structures on the noisy and hysteretic examples. The formulation is
supported by proofs that the identified model and both objectives depend only on the executed set of monomials, that
the arithmetic cost is modular and strictly monotone for every weight and encodes the number of terms and the total
degree without loss under the bit-level weight, and that constant cost factors have no effect; the lexicographic mode
was shown to reproduce the single-objective search except for ties in the error, and the Pareto-based scheme to
preserve its population front under a capacity condition, both of which were confirmed generation by generation.
These results extend Green-Box identification, previously embedded in greedy orthogonal-least-squares selection, to
evolutionary structure search, and they show that an operation-level cost changes the identified structures only when
the selection scheme treats it as an objective rather than as a tie-breaker, and that the weight of the cost is a
property of the executing platform that should be calibrated before identification. The lexicographic mode improved
on the single-objective search mainly through its final non-dominated filtering. The conclusions are limited to small
examples and to the model class that the package can currently evaluate, which excludes individuals without any
backshift operator and products of delayed outputs and inputs with delays greater than one, and the single-model rule
used only identification error, so that on noisy data it preferred structures with spurious terms over the true one
present in the same set, and on the hysteretic example it selected cheaper models with a slightly, though not
significantly, higher validation error. The execution measurements concern interpreted code on a single laptop
processor in one session, and the package's own simulator, which evaluates duplicated genes, realised a smaller part of
the saving. Future work will remove these package restrictions, combine the selection rule with validation-based or
information criteria, extend the comparison to established benchmarks and to the orthogonal-least-squares Green-Box
formulation, and calibrate the weight on the embedded processors where the identified models are intended to run.
